\section{Event-triggered realization}
\label{sec:events}

\subsection{Events, handlers and inversion of control}
\label{sec:events-ioc}

A graph and an event loop can run the same machine with different mechanisms for invoking
application code. In the graph, a routing function selects the next node, and the scheduler runs it
in the next superstep (\cref{sec:boundary-scheduling}). In an event loop, a computation emits an
event, and the runtime invokes the code registered to handle it. The following four definitions
describe this arrangement. Their wording comes from one source~\citep[Event Handling; Understanding the
Event Loop]{cornell-cs2110-lec25}.

\begin{definition}[Event]
\label{def:event}
``An event is any occurrence that could necessitate a change in the program's state.''
\end{definition}

\begin{definition}[Handler]
\label{def:handler}
Handlers are ``pieces of code that understand how to respond to a particular event''; for events to
reach them, they ``must register with the event loop.''
\end{definition}

\begin{definition}[Event loop]
\label{def:event-loop}
``An event loop is a piece of code that repeatedly checks whether a new event has occurred and
dispatches any detected events to code that can handle them.''
\end{definition}

\begin{definition}[Inversion of control]
\label{def:inversion-of-control}
``Inversion of control is used to describe the paradigm where we write custom subroutines but allow
an external entity, such as a framework, to control when these subroutines are executed.''
\end{definition}

For an agent, an event represents either an occurrence, such as \emph{the plan is ready} or
\emph{the checks failed}, or a request, such as \emph{run these tools}. Both are values. A tool
request event records a request to a handler. It does not establish that the handler ran. The same
distinction applies to a model's tool request in \cref{sec:boundary-artifacts}: requesting a file
read does not establish that the read occurred.

Inversion of control separates two responsibilities. Application code defines what each handler
does with its event, which fields of the record it writes, and which event it emits on completion.
The framework determines which handler runs, when it runs, and the order in which it processes
pending events. \textbf{Application code defines the response to an event; the framework controls
its invocation.} The graph realization makes a similar division: we write node bodies and routing
functions, and the scheduler runs them. The two realizations express the choice of the next
computation in different places. A graph uses a routing function on an edge. An event handler
chooses the type of event to emit, and the registrations map that type to a handler.

The registrations determine which event types have handlers. If the loop does not check its
registrations, it can remove an event from the queue even though no handler accepts its type.
Nothing then runs, although no component has decided to stop. Whether the framework reports this
as an error depends on its implementation. We treat it as a defect rather than a terminal
outcome.

\subsection{An event-triggered version of the machine}
\label{sec:events-machine}

Changing the realization changes when the runtime takes a transition. The set of transitions stays
the same. A state-machine diagram does not specify when a reaction occurs. In an event-triggered
model, a reaction occurs when an input is present; in a time-triggered model, it occurs on the tick of an
external clock~\citep[slides~23--25]{seshia2015discrete}. The planning agent's machine
(\cref{fig:efsm-planner}) retains its locations, guards and set actions under either choice.
Our graph realization uses no clock: each superstep begins when the preceding one finishes. The
event realization reacts when an event is present and is therefore event-triggered under this
definition.

The choice can still make a transition unreachable. A guard requiring every triggering input to be
\emph{absent} cannot fire when reactions require a triggering input to be present. Each transition
of the planning agent becomes enabled when the preceding computation completes, so each can wait
for a completion event. A transition such as \emph{no reply
arrived within thirty seconds} needs a timer to emit an event when the interval ends. Without the
timer, the runtime never takes this transition even though it appears in the diagram.

The event realization requires a larger configuration. In the graph realization, a configuration
contains a control location and a valuation of the variables, $(q, v)$ (\cref{def:configuration}).
In the event realization it is a triple
\[
  (E,\ \mathcal{H},\ v),
\]
where $E$ is the sequence of pending events, $\mathcal{H}$ is the set of running handlers, and $v$
is the record. The runtime maintains $E$ and $\mathcal{H}$. The record retains the steps' data and
contains neither the queue nor the set of running handlers. In a serial realization, at most one
event is in flight. During active execution, either one event is pending and no handler runs, or
one handler runs and no event is pending. We can then recover the location $q$ from the handler
registered for the pending event or from the running handler. A realization that permits two
handlers to run at once can have configurations with no single location, and the machine in
\cref{fig:efsm-planner} no longer describes such configurations. Our demo's validation step is
such a case (\cref{sec:events-demo}).

Four rules preserve the machine's behaviour when an event loop realizes it.

\bi

\item \textbf{Apply once, then enable.} Merge a handler's contribution into the record exactly
once before making its emitted event available for dispatch. The graph's superstep model provides
this ordering (\cref{def:superstep}). An event loop provides it only if the handler writes before
it emits and the runtime does not deliver the same event twice. If a handler emits first, its successor can
read the record before the contribution appears. Delivering an event twice appends the same
observations twice.

\item \textbf{Distinguish idle from success.} An empty queue with no running handler means that
nothing is scheduled. It does not establish how the run ended or whether termination was
intentional. Only a terminal event carrying a status establishes termination, as a terminal record
does in the trace (\cref{sec:control-termination}).

\item \textbf{Stop dispatch at termination.} After dispatching the terminal event, the runtime
must start no further handler, even if an event remains pending. In a serial realization, no event
can be pending at that point. A realization with concurrent handlers must discard or refuse
the remaining events.

\item \textbf{Separate triggering events from log events.} A triggering event is an input to the
dispatcher that causes a handler to run. A log event records an occurrence in the trace and
triggers no computation. An event that requests tools triggers the tools handler; a trace
record of the same request documents that the handler was invoked, and triggers nothing.

\ei

The third and fourth rules both apply at termination. The terminal event ends dispatch, and a
terminal record in the trace records its status. If the process fails before writing that record,
dispatch can end without a terminal record. The interpretations in \cref{tab:control-outcomes}
still apply.

\subsection{Two frameworks}
\label{sec:events-frameworks}

Two frameworks provide event dispatch through different APIs.

\bi

\item \textbf{LlamaIndex Workflows.} Dispatch is by event type. A step is a method decorated with
\texttt{@step}, and ``A step receives an event, does some work, and returns another event. That
returned event triggers the next step whose type annotation accepts it''~\citep{llamaindex-workflows}.
Events are user-defined Pydantic classes, and the framework supplies two of its own:
\texttt{StartEvent}, whose fields are the keyword arguments passed to \texttt{run()}, and
\texttt{StopEvent}, whose return ends the workflow with the value passed as its \texttt{result}.
The step signatures declare the event types each step consumes and produces. Before a run, the
framework validates the event graph, checking, among other conditions, that it contains a start
and a stop event and that every produced event has a consumer. A store on the workflow's
\texttt{Context} holds per-run state. Steps read and write it through \texttt{ctx.store}, using
\texttt{edit\_state()} for updates that must not interleave with another step's
updates~\citep{llamaindex-workflows-state}. Our demo is written in this
framework (\cref{sec:events-demo}).

\item \textbf{CrewAI Flows.} Dispatch uses method completion and route labels. A method decorated
with \texttt{@start()} is an entry point. A method decorated with \texttt{@listen(m)} runs when
method \texttt{m} completes and receives its output. A method decorated with \texttt{@router(m)}
runs after \texttt{m} and returns a string label. Methods decorated with \texttt{@listen("label")}
run when the router returns that label. The operators \texttt{or\_} and \texttt{and\_} combine
triggers: a listener runs when any or all of several methods have produced output, respectively.
The flow's \texttt{state} attribute holds shared state as a dictionary or a Pydantic model. There is
no stop event: ``the final output is determined by the last method that completes''~\citep{crewai-flows}.
Translated to CrewAI, each of our event types would become a router label, and the outcome
$\sigma$ would have to be written to \texttt{state} by our code, because the end of the flow does
not carry it.

\ei

This difference affects how we distinguish idle from success. LlamaIndex has an explicit terminal
event. Its pre-run validation checks the event graph declared by step signatures for start and
stop events and missing event connections. These checks do not establish that a particular run
will reach \texttt{StopEvent}. CrewAI defines no terminal event. Our reading of its documentation
is that a flow ends when no method remains to run, corresponding to an empty queue. The outcome
must therefore be in the record before the last method completes.

Neither framework establishes that an event realization preserves the machine's behaviour. Comparing it
with the graph realization requires a controlled design. We compare the realizations after each
computation completes and its contribution is applied. These logical boundaries define the points
of comparison; internal handler steps and wall-clock ordering do not. We script the model's replies
and control the tool results so that both realizations receive the same inputs in the same order.
A faithful translation preserves four properties at each boundary: i) the observations appended
to the record, ii) the effects on the workspace, iii) the results of the acceptance checks, and
iv) the terminal status and its meaning. The comparison is designed to isolate the scheduling
mechanism: a graph superstep or an event dispatch.

Agreement under scripted inputs supports a limited conclusion: the two realizations implement the
same transition rules for the scripted replies. It does not establish equivalent behaviour with a
live model. Model replies depend on the messages each realization constructs, so changing how a
handler formats messages changes the input distribution even when the transition rules stay the
same. The comparison also excludes concurrency.

\subsection{The demo: the planning agent as a LlamaIndex workflow}
\label{sec:events-demo}

Our event-driven agent, \demoref{cyberbird/event\_driven}, runs the planning agent's machine as a
LlamaIndex workflow. It keeps the machine's locations, guards and limits: a budget of forty model
calls per alert, checked before each call; a stall after three identical results; and one plan
revision. The run's data sit in a store on the workflow's context. A \demoref{RunState} holds the
run's bookkeeping, and one \demoref{AlertState} per alert holds that alert's plan, messages,
observations, calls and usage. Two of its steps are still stubs: \demoref{submit} runs no
mechanical checks yet, and each validator accepts after a fixed delay without calling a model. No
run of this agent has been recorded.

\begin{table*}[tbp]
  \centering
  \begin{adjustbox}{width=\linewidth}\begin{tabularx}{\linewidth}{@{}L{0.26\linewidth}L{0.10\linewidth}XL{0.10\linewidth}@{}}
    \toprule
    \textbf{Events the step accepts} & \textbf{Step} & \textbf{Events the step emits} & \textbf{Decided by} \\
    \midrule
    \texttt{RunRequested} & \demoref{start} & \texttt{ScanRequested} & runtime \\
    \texttt{ScanRequested} & \demoref{scan} & one \texttt{AlertFound} per alert, or \texttt{StopEvent} & runtime \\
    \texttt{AlertFound}, \texttt{AlertDone} & \demoref{gate} & \texttt{AlertStarted}, one alert at a time, or \texttt{StopEvent} & runtime \\
    \midrule
    \texttt{AlertStarted}, \texttt{Stalled} & \demoref{plan} & \texttt{PlanCreated}, \texttt{AlertDone} & model \\
    \texttt{PlanCreated}, \texttt{Observed}, \texttt{ChecksFailed} & \demoref{controller} & \texttt{ToolsRequested}, \texttt{Submitted}, \texttt{AlertDone} & model \\
    \texttt{ToolsRequested} & \demoref{tools} & \texttt{Observed}, \texttt{Stalled}, \texttt{AlertDone} & runtime \\
    \texttt{Submitted} & \demoref{submit} & \texttt{ChecksFailed}, or $V$ $\times$ \texttt{ValidateRequested} & runtime \\
    \texttt{ValidateRequested} & \demoref{validate} & \texttt{Verdict} & model \\
    $V$ $\times$ \texttt{Verdict} & \demoref{judge} & \texttt{AlertDone} & runtime \\
    \bottomrule
  \end{tabularx}\end{adjustbox}
  \caption{The steps of the event-driven agent, from the wiring table in its \demoref{workflow.py}.
  Each step's signature names the events it accepts, and its return annotation names the events it
  emits. The rows below the rule are the machine of \cref{fig:efsm-planner}; the rows above it are
  new. $V$ is the number of validators, three by default.}
  \label{tab:events-steps}
\end{table*}

The six rows below the rule in \cref{tab:events-steps} are the machine's locations. No step names
another: the wiring is the set of step signatures, and the framework dispatches each event to the
step whose signature accepts its type. The guards that the graph evaluated in routing functions
now sit inside the steps and choose which event to return.

\bi

\item \textbf{Admission.} \demoref{plan} and \demoref{controller} first reserve one model call for
the alert. When the alert has spent its budget, the step returns \texttt{AlertDone} with status
\demoref{budget\_exhausted} and makes no call.

\item \textbf{Acting.} The controller returns \texttt{ToolsRequested} when the model's reply
contains tool calls, and \texttt{Submitted} when it contains none.

\item \textbf{Revising.} \demoref{tools} returns \texttt{Observed} while the work moves. On the
first stall it returns \texttt{Stalled}, which the plan step accepts; that step writes a new plan
and moves \demoref{replanned\_at}, the handoff and window of \cref{sec:state-lifetimes}. A second
stall returns \texttt{AlertDone} with status \demoref{no\_progress}.

\ei

The event realization also does three things the graph agent does not. First, one run handles
many alerts. \demoref{scan} runs Bandit once and emits one \texttt{AlertFound} per alert, and
\demoref{gate} releases one alert at a time: it queues each alert it receives and starts the next
one only after the current alert's \texttt{AlertDone} arrives. Every event about an alert carries
its \demoref{alert\_id}, which correlates the event with that alert's entry in the store.

Second, validation fans out and back in. \demoref{submit} sends one \texttt{ValidateRequested} per
validation lens (correctness, behaviour and scope) with \demoref{ctx.send\_event}, because a step
can return only one event. The validate step runs with up to three workers, so the three
validators run at the same time. This is the configuration of \cref{sec:events-machine} with more
than one running handler and no single location. \demoref{judge} holds each \texttt{Verdict}
until all $V$ have arrived, in any order, and accepts the alert only when every validator accepted.
It collects verdicts without keying them by alert, which is safe only because the gate admits one
alert at a time. \textbf{The gate serializes alerts; it does not serialize validators.}

Third, the run has two kinds of ending. An alert ends with an \texttt{AlertDone} carrying one of
four statuses: \demoref{accepted}, \demoref{rejected}, \demoref{no\_progress} or
\demoref{budget\_exhausted}. That event returns to the gate, which records the status. The run
ends when the gate has a result for every alert the scan found, and it returns a
\texttt{StopEvent} carrying the results. A scan that finds no alert ends the run at once with an
empty result.

\begin{figure*}[tbp]
  \centering
  \begin{adjustbox}{width=\linewidth}% The event-driven agent (fig:event-loop): cyberbird/event_driven/workflow.py.
% Each step is drawn with the event types its signature accepts; every event a step returns or
% sends goes back to the queue. Styles from the lecture_notes.tex preamble.
\begin{tikzpicture}[
  rt/.style={gnode, minimum width=20mm},
  hm/.style={mnode, minimum width=21mm},
  hr/.style={rnode, minimum width=21mm},
  ev/.style={lab, font=\ttfamily\scriptsize, align=left},
  every node/.append style={text=ibmfg}
]
  % runtime components
  \node[rt] (queue) at (0,0) {event queue\\(pending events)};
  \node[rt] (disp) at (3.9,0) {dispatcher};

  % steps, top to bottom in the order an alert meets them
  \node[hr] (start) at (11.4,5.6) {start};
  \node[hr] (scan) at (11.4,4.2) {scan};
  \node[hr] (gate) at (11.4,2.8) {gate};
  \node[hm] (plan) at (11.4,1.4) {plan};
  \node[hm] (ctl) at (11.4,0) {controller};
  \node[hr] (tools) at (11.4,-1.4) {tools};
  \node[hr] (submit) at (11.4,-2.8) {submit};
  \node[hm] (val) at (11.4,-4.2) {validate $\times V$};
  \node[hr] (judge) at (11.4,-5.6) {judge};

  % start and end of the run
  \node[tnode] (req) at (-3.6,0) {run request};
  \node[tnode, minimum width=14mm] (end) at (3.9,-4.2) {END};

  \draw[flow] (req) -- node[lab, above] {filters} (queue);
  \draw[flow] (queue) -- node[lab, above] {next\\event} (disp);

  % dispatch: each arrow carries the event types its step accepts
  \draw[thick, draw=ibmfg!80] (disp.east) -- (6.2,0);
  \draw[thick, draw=ibmfg!80] (6.2,5.6) -- (6.2,-5.6);
  \draw[flow] (6.2,5.6) -- node[ev, above] {RunRequested} (start.west);
  \draw[flow] (6.2,4.2) -- node[ev, above] {ScanRequested} (scan.west);
  \draw[flow] (6.2,2.8) -- node[ev, above] {AlertFound, AlertDone} (gate.west);
  \draw[flow] (6.2,1.4) -- node[ev, above] {AlertStarted, Stalled} (plan.west);
  \draw[flow] (6.2,0) -- node[ev, above] {PlanCreated, Observed,\\ChecksFailed} (ctl.west);
  \draw[flow] (6.2,-1.4) -- node[ev, above] {ToolsRequested} (tools.west);
  \draw[flow] (6.2,-2.8) -- node[ev, above] {Submitted} (submit.west);
  \draw[flow] (6.2,-4.2) -- node[ev, above] {ValidateRequested} (val.west);
  \draw[flow] (6.2,-5.6) -- node[ev, above] {Verdict} (judge.west);

  % the stop event ends the run
  \draw[flow] (disp) -- node[lab, left] {StopEvent\\with the results} (end);

  % every event a step emits returns to the queue
  \foreach \h in {start, scan, gate, plan, ctl, tools, submit, val, judge} {
    \draw[dashed, thick, draw=ibmfg!80] (\h.east) -- (13.4,0 |- \h.east);
  }
  \draw[back] (13.4,-5.6) -- (13.4,6.7) -| node[lab, above, pos=.25] {emitted events} (queue.north);
\end{tikzpicture}
\end{adjustbox}
  \caption{The event-driven agent of \cref{tab:events-steps}. White boxes are runtime components.
  Steps use the colours of \cref{fig:efsm-planner}: blue steps call a model (the validators will,
  once they are no longer stubs), and grey steps run tools, checks or bookkeeping. The ellipses
  are the run request that starts the run and \texttt{END}. Each solid arrow from the dispatcher
  names the event types its step accepts. Every event a step emits returns to the queue along the
  dashed path. $V$ validators run at the same time.}
  \label{fig:event-loop}
\end{figure*}

The code keeps the four rules of \cref{sec:events-machine} in these forms.

\bi

\item \textbf{Apply once, then enable.} Every write goes through \demoref{ctx.store.edit\_state()},
which locks the store, and each step writes its data before it emits the event that lets another
step read them. The budget check reserves its call under the same lock, so two steps admitted at
the same moment cannot both take the last call.

\item \textbf{Distinguish idle from success.} Only the gate's \texttt{StopEvent} ends the run. Before the
first run, the framework also checks that every event a step emits has a step that accepts it,
which is why \demoref{submit} lists \texttt{ValidateRequested} in its return annotation although it
sends those events rather than returning them.

\item \textbf{Stop dispatch at termination.} The gate emits \texttt{StopEvent} only when the count of
results equals the count of alerts, so no alert is still in progress when the run ends.

\item \textbf{Separate triggering events from log events.} This agent writes no trace. Its events trigger
steps, and no event records one. Our graph agents write a trace alongside their events
(\cref{sec:control-termination}); the event-driven agent does not yet have one.

\ei

\subsection{Reading}

\citet{cornell-cs2110-lec25}, the sections Event Handling and Understanding the Event
Loop, for events, handlers, the event loop and inversion of control; and \citet{seshia2015discrete},
slides 23--25, on event-triggered and time-triggered reactions.

For reference while writing HW1: \citet{llamaindex-workflows}, the sections on entry points, exit
points and validation, with \citet{llamaindex-workflows-state} for the per-run store; and
\citet{crewai-flows}, the sections on \texttt{@start()}, \texttt{@listen()}, retrieving the final
output, conditional logic and the router.

\paragraph{Looking ahead.} In an event loop, a registered handler runs when its event arrives, and
a handler can emit any number of events. The loop itself does not decide whether the work is
permitted or assess its cost. The next section distinguishes admission, which decides whether a
model or tool call may start, from accounting, which records its cost after execution. It also
considers how to measure a realization's cost under the controlled conditions used here to compare
behaviour.
